% --- 문제풀이
\begin{frame}{문제풀이}
\vfill
\begin{enumerate}
\item 1차원 점: 빨강 $x = 0.5, 1.5$, 파랑 $x = -1.5, -0.5$. (a) 최대 마진 경계와 마진 $m$ 은? \; (b) 서포트 벡터는?\vskip0.8em
\item 경계 $0.6x_1 + 0.8x_2 - 1 = 0$. (a) $\lVert\mathbf{w}\rVert$ 는? \; (b) 점 $(1, 1)$ 에서 경계까지 거리는?\vskip0.8em
\item 학습 정확도 100\%, 검증 정확도 70\% 인 RBF SVM. sklearn 의 $C$ 와 $\gamma$ 를 어느 쪽으로 바꿔 보겠나?\vskip0.8em
\item $\mathbf{x} = (2, 0)$, $\mathbf{x}' = (1, 3)$ 일 때 $(\mathbf{x}\cdot\mathbf{x}')^2$ 와 $\phi(\mathbf{x})\cdot\phi(\mathbf{x}')$ 를 각각 계산해 같은지 확인하라.\vskip0.8em
\item 클래스 5개를 SVM 으로 분류한다. 일대일 · 일대다는 각각 SVM 이 몇 개 필요한가?
\end{enumerate}
\vfill
\note{1번 답. 가장 가까운 두 점이 마이너스 0.5와 0.5이니 경계는 가운데 0, 마진 m은 0.5, 길 폭은 1입니다. 서포트 벡터는 마이너스 0.5와 0.5입니다. 2번 답. w의 길이는 루트 0.36 더하기 0.64, 1입니다. 그러니 값이 곧 거리라서, 0.6 더하기 0.8 빼기 1, 거리 0.4입니다. 3번 답. 과대적합이니 C와 감마를 둘 다 줄여 봅니다. 4번 답. 점곱이 2라서 제곱하면 4입니다. 파이 x는 (4, 0, 0), 파이 x 프라임은 (1, 3루트2, 9)이니 점곱은 4, 같습니다. 5번 답. 일대일은 5 곱하기 4 나누기 2, 10개, 일대다는 5개입니다.}
\end{frame}


\begin{frame}{문제풀이: 학습 한 걸음 손으로}
\vfill
\centering
{\large 빨강 ($y=+1$): $A(1.5, 0)$, \; $B(0, 1.5)$ \qquad 파랑 ($y=-1$): $C(0, -1.5)$}\\[0.5em]
{\large 지금 선: $\mathbf{w} = (0, 1)$, \; $b = 0.5$ \quad ($x_2 = -0.5$), \quad $\eta = 0.5$}\\[0.9em]
\raggedright
\begin{enumerate}
\item[(a)] 세 점의 값 $y_i(\mathbf{w}\cdot\mathbf{x}_i + b)$ 와 $m$ 은? 어느 점이 미나?\vskip0.5em
\item[(b)] $\mathbf{w}' = \mathbf{w} + \eta\, y_k \mathbf{x}_k$, \; $b' = b + \eta\, y_k$ 를 계산하라.\vskip0.5em
\item[(c)] $\lVert\mathbf{w}'\rVert$ 는? 길이 1 로 맞추면 $\mathbf{w}, b$ 는?\vskip0.5em
\item[(d)] 새 선에서 세 점의 값과 $m$ 은? 다음 걸음엔 어느 점이 미나?
\end{enumerate}
\vfill
\note{(a) A는 0 더하기 0.5, 0.5. B는 1.5 더하기 0.5, 2. C는 파랑이라 마이너스 곱하기 (마이너스 1.5 더하기 0.5), 1. 가장 작은 값은 A의 0.5이니 m은 0.5이고 A가 밉니다. (b) A는 빨강, y는 플러스 1. w 프라임은 (0, 1) 더하기 0.5 곱하기 (1.5, 0), (0.75, 1). b 프라임은 0.5 더하기 0.5, 1. (c) w 프라임의 길이는 루트 0.5625 더하기 1, 루트 1.5625, 1.25입니다. 1을 넘으니 둘 다 1.25로 나눠 w는 (0.6, 0.8), b는 0.8. 선은 0.6 x1 더하기 0.8 x2 더하기 0.8 = 0입니다. (d) A는 0.9 더하기 0.8, 1.7. B는 1.2 더하기 0.8, 2. C는 마이너스 곱하기 (마이너스 1.2 더하기 0.8), 0.4. m은 0.4로 오히려 줄었습니다. A에서 너무 멀리 밀어서 C 쪽으로 넘어간 것입니다. 다음 걸음엔 C가 밉니다. 이렇게 한 걸음에 m이 꼭 늘지는 않고, 가장 가까운 점이 번갈아 밀면서 맞춰 갑니다. 그래서 학습률을 걸음이 갈수록 줄입니다.}
\end{frame}

